\newpage
\phantomsection
\label{prf:scalar-mult-negative}
\LRAProofFor{thm:scalar-mult-negative}

\begin{remark*}[Return]
\hyperref[thm:scalar-mult-negative]{Return to Theorem}
\end{remark*}

\ProofVaultURL{https://github.com/wsollers/lra-proof-vault/tree/master/volume-iii/book-analysis-i/bounding/thm-scalar-mult-negative}

\begin{theorem*}[Scalar Multiplication by a Negative Scalar]
Let $A\subseteq\mathbb{R}$ be nonempty and bounded below, and let $k<0$. Then
\[
  \sup(kA)=k\inf A,
  \qquad
  kA=\{ka:a\in A\}.
\]
\end{theorem*}

\begin{proof}[Professional Standard Proof]
\LRAProofBodyStart
Let $N=-k$. Since $k<0$, we have $N>0$. Also
\[
  kA=\{ka:a\in A\}=\{N(-a):a\in A\}=N(-A).
\]
Because $A$ is bounded below, the set $-A$ is bounded above. By scalar
multiplication by a positive scalar,
\[
  \sup(N(-A))=N\sup(-A).
\]
By the theorem that negation exchanges supremum and infimum,
\[
  \sup(-A)=-\inf A.
\]
Therefore
\[
  \sup(kA)=\sup(N(-A))=N\sup(-A)=N(-\inf A)=-N\inf A.
\]
Since $N=-k$, we have $-N=k$, and hence $\sup(kA)=k\inf A$.
\end{proof}

\begin{proof}[Detailed Learning Proof]
\LRAProofBodyStart
We reduce the negative-scalar case to two results that are already known:
positive scalar multiplication of suprema, and the exchange of supremum and
infimum under negation.

Let $N=-k$. The hypothesis $k<0$ implies $N>0$. For every $a\in A$,
\[
  ka=(-N)a=N(-a).
\]
Therefore the image of $A$ under multiplication by $k$ is the same set as the
image of $-A$ under multiplication by the positive scalar $N$:
\[
  kA=N(-A).
\]

Since $A$ is bounded below, negating all of its elements produces a set
$-A=\{-a:a\in A\}$ that is bounded above. The positive-scalar theorem applies
to the nonempty bounded-above set $-A$ and the positive scalar $N$, so
\[
  \sup(N(-A))=N\sup(-A).
\]
The negation theorem gives
\[
  \sup(-A)=-\inf A.
\]
Substituting this into the preceding equality yields
\[
  \sup(N(-A))=N(-\inf A).
\]

Finally, because $kA=N(-A)$, the two sets have the same supremum. Hence
\[
  \sup(kA)=\sup(N(-A))=N(-\inf A)=-N\inf A.
\]
Since $N=-k$, we have $-N=k$. Thus
\[
  \sup(kA)=k\inf A,
\]
as required.
\end{proof}

\begin{remark*}[Proof structure]
Convert the negative multiplier $k$ into a positive multiplier $N=-k$ applied
to the negated set $-A$. Then combine the positive-scalar supremum formula
with the theorem that negation exchanges supremum and infimum.
\end{remark*}

\begin{dependencies}
\begin{itemize}
  \item \hyperref[thm:scalar-mult-positive]{Scalar Multiplication by a Positive Scalar}.
  \item \hyperref[thm:negation-exchange-sup-inf]{Negation Exchanges Supremum and Infimum}.
  \item \hyperref[def:supremum]{Supremum}.
  \item \hyperref[def:infimum]{Infimum}.
\end{itemize}
\end{dependencies}

\clearpage
